\documentclass[12pt,a4paper]{article}

% Encoding and fonts — LuaLaTeX
\usepackage{fontspec}
\usepackage{unicode-math}
\setmainfont{FreeSerif}
\setmathfont{Latin Modern Math}
\newfontfamily\igprimfont{FreeSerif}
\setmonofont[Scale=0.85]{DejaVu Sans Mono}

% Page layout
\usepackage[top=1in, bottom=1in, left=1in, right=1in]{geometry}
\setlength{\headheight}{15pt}

% Typography
\usepackage{microtype}
\usepackage{parskip}

% Language
\usepackage[english]{babel}

% Headers and footers
\usepackage{fancyhdr}
\pagestyle{fancy}
\fancyhf{}
\fancyhead[L]{\leftmark}
\fancyhead[R]{\thepage}
\fancyfoot[C]{\thepage}
\renewcommand{\headrulewidth}{0.4pt}
\renewcommand{\footrulewidth}{0pt}

% Section styling
\usepackage{titlesec}
\titleformat{\section}{\Large\bfseries}{\thesection}{1em}{}
\titleformat{\subsection}{\large\bfseries}{\thesubsection}{1em}{}
\titlespacing*{\section}{0pt}{2.5ex plus 1ex minus .2ex}{1.5ex plus .2ex}
\titlespacing*{\subsection}{0pt}{2ex plus 1ex minus .2ex}{1ex plus .2ex}

% Lists
\usepackage{enumitem}
\setlist[itemize]{topsep=0.5ex, itemsep=0.25ex, parsep=0pt}
\setlist[enumerate]{topsep=0.5ex, itemsep=0.25ex, parsep=0pt}

% Tables and math
\usepackage{graphicx}
\usepackage[table]{xcolor}
\usepackage{booktabs}
\usepackage{array}
\usepackage{tabularx}
\usepackage{amsmath}
\usepackage{amsthm}
\usepackage{calc}
\usepackage[normalem]{ulem}
\renewcommand{\arraystretch}{1.4}

% Cross-references
\usepackage{hyperref}
\usepackage{cleveref}
\hypersetup{colorlinks=true, linkcolor=blue, urlcolor=cyan}

% Code listings
\usepackage{listings}
\definecolor{codebg}{RGB}{248,248,248}
\definecolor{codeframe}{RGB}{200,200,200}
\lstset{
    basicstyle=\ttfamily\small,
    breaklines=true,
    frame=single,
    framerule=0.5pt,
    rulecolor=\color{codeframe},
    backgroundcolor=\color{codebg},
    keepspaces=true,
    numbers=left,
    numberstyle=\tiny\color{gray},
    numbersep=8pt,
    showstringspaces=false,
    tabsize=4,
    xleftmargin=1em,
    xrightmargin=1em,
    aboveskip=1em,
    belowskip=1em,
    literate={_}{{\textunderscore}}1
             {^}{{\textasciicircum}}1
}

% Boxes
\usepackage{tcolorbox}
\tcbuselibrary{skins,breakable}
\newtcolorbox{imscriptionbox}{
    enhanced,
    colback=blue!5,
    colframe=blue!75!black,
    fontupper=\large,
    center upper,
    boxrule=1pt,
    arc=4pt,
}
\newtcolorbox{admonitionbox}[2][]{
    enhanced, breakable,
    colback=#2!5, colframe=#2!80!black,
    fonttitle=\bfseries, title=#1,
    left=2mm, right=2mm,
}

% Theorem environments
\newtheorem{theorem}{Theorem}
\newtheorem{lemma}[theorem]{Lemma}
\newtheorem{corollary}[theorem]{Corollary}
\theoremstyle{definition}
\newtheorem{definition}{Definition}
\theoremstyle{remark}
\newtheorem{remark}{Remark}

\newcommand{\tightlist}{\setlength{\itemsep}{0pt}\setlength{\parskip}{0pt}}

\begin{document}

\title{The Frobenius Loop:\\
A Permanently Paradoxical Fixed Point\\
in the Paraconsistent Universal Engine}
\author{umpolungfish}
\date{2026-05-21}
\maketitle

\begin{abstract}
We exhibit a three-instruction circular program on the TriPhaseRegister
machine---the Universal Engine extended with Belnap FOUR belief
tracking---that runs indefinitely with all registers permanently in the
\emph{Both} (B) state.  The loop invariant is a computational realization of
the Frobenius identity $\mu \circ \delta = \mathrm{id}$: FFUSE applied to two
B-state outputs of FSPLIT returns B, preventing collapse to any classical
truth value.  We prove that the paradox count after $n$ complete cycles is
exactly $P(n) = 4n$, and confirm this algebraic prediction empirically at
$1.9 \times 10^{10}$ steps of continuous execution.
\end{abstract}

\section{The Machine}

\begin{definition}[TriPhaseRegister]
A \emph{TriPhaseRegister} is a pair $(\mathrm{flux}, \mathrm{val})$ where
$\mathrm{flux} \in \{00, 01, 10, 11\}$ and $\mathrm{val} \in \{\bot,
\texttt{FIXED}\}$.  The flux encodes a two-bit state:
\[
  00 = \text{Void},\quad
  01 = \text{True},\quad
  10 = \text{False},\quad
  11 = \text{Both}.
\]
Each register also carries a non-negative integer \emph{paradox count}
$\pi \in \mathbb{N}$, incremented by every engagement event.
\end{definition}

\begin{definition}[Belnap FOUR]
The Belnap four-valued lattice $\mathbf{B}_4 = \{N, T, F, B\}$ is equipped
with two partial orders.  The \emph{truth order} has $N, T, F < B$ and $N <
T, F$.  The \emph{information order} has $N < T$ and $N < F$, with $T, F$
incomparable and both below $B$.  The \emph{Belnap join} $\vee$ is the join
in the information order:
\[
  N \vee x = x,\quad T \vee F = B,\quad B \vee x = B,\quad x \vee x = x.
\]
In particular $B \vee B = B$ (idempotency of paradox).
\end{definition}

The \emph{Paraconsistent Universal Engine} (\texttt{ParaEngine}) is the
TriPhaseRegister machine augmented with a belief map
$\beta : \mathrm{Reg} \to \mathbf{B}_4$ kept in sync with flux:
$\beta(r) = N, T, F, B$ iff $r.\mathrm{flux} = 00, 01, 10, 11$
respectively.  The relevant IMASM opcodes are:

\begin{itemize}
  \item \texttt{ENGAGR \%r$_i$}: sets $r_i.\mathrm{flux} \leftarrow 11$,
        $\beta(r_i) \leftarrow B$, $r_i.\pi \mathrel{+}= 1$.
  \item \texttt{FSPLIT \%r$_i$ \%r$_j$ \%r$_k$}: applies \texttt{ENGAGR}
        to each of $r_i, r_j, r_k$ in turn.  Total: three engagement events,
        $\pi$ incremented for each.
  \item \texttt{FFUSE \%r$_i$ \%r$_j$ \%r$_k$}: sets
        $\beta(r_k) \leftarrow \beta(r_i) \vee \beta(r_j)$ (Belnap join).
        No engagement events; $\pi$ unchanged.
\end{itemize}

FSPLIT is the Frobenius co-multiplication $\delta$; FFUSE is the
multiplication $\mu$.

\section{The Kernel}

The \emph{Frobenius loop kernel} $\mathcal{K}$ is the following
three-instruction program:

\begin{lstlisting}[caption={Frobenius loop kernel},label={lst:kernel}]
K[0]  ENGAGR %r0          ; seed: force B on the root register
K[1]  FSPLIT %r0 %r1 %r2  ; delta: co-multiply -- two B-children
K[2]  FFUSE  %r1 %r2 %r0  ; mu: multiply -- fold back into root
\end{lstlisting}

The program counter is circular: after executing \texttt{K[2]} the engine
sets $\mathrm{pc} \leftarrow 0$ and continues.  This is not an imposed
wrapper; the \texttt{step()} method of the Universal Engine resets to zero
at program end by design (\emph{bootstrap loop closure}).

\section{The Fixed Point Theorem}

\begin{theorem}[Permanence of B-state]
\label{thm:permanent-b}
Let the engine start in any state.  After the first complete cycle of
$\mathcal{K}$, we have $\beta(r_0) = \beta(r_1) = \beta(r_2) = B$.  This
holds for every subsequent cycle.
\end{theorem}

\begin{proof}
Track beliefs through one cycle, starting from an arbitrary initial state
$\beta(r_0) = s_0 \in \mathbf{B}_4$.

\begin{enumerate}
  \item \texttt{ENGAGR \%r0}: $\beta(r_0) \leftarrow B$ (forced,
        regardless of $s_0$).
  \item \texttt{FSPLIT \%r0 \%r1 \%r2}: $\beta(r_0) \leftarrow B$,
        $\beta(r_1) \leftarrow B$, $\beta(r_2) \leftarrow B$ (each register
        is engaged independently).
  \item \texttt{FFUSE \%r1 \%r2 \%r0}: $\beta(r_0) \leftarrow \beta(r_1)
        \vee \beta(r_2) = B \vee B = B$ (Belnap join idempotency).
\end{enumerate}

After cycle 1: $\beta(r_0) = \beta(r_1) = \beta(r_2) = B$.  For any later
cycle, \texttt{ENGAGR} at \texttt{K[0]} forces $B$ on $r_0$ (already $B$),
\texttt{FSPLIT} forces $B$ on all three (already $B$), and \texttt{FFUSE}
computes $B \vee B = B$.  By induction the belief triple $(B, B, B)$ is a
fixed point of $\mathcal{K}$.
\end{proof}

\begin{remark}
The Frobenius identity $\mu \circ \delta = \mathrm{id}$ underlies step (3):
FFUSE undoes the split performed by FSPLIT, mapping the B-pair back to the
original B.  Without idempotency of $\vee$---i.e., in a lattice where $B
\vee B \neq B$---this fixed point need not exist.  The Belnap lattice is
precisely the setting in which paraconsistent contradiction is
\emph{self-sustaining} rather than explosive.
\end{remark}

\section{The Paradox Growth Theorem}

Let $P(n)$ denote the total paradox count
$\sum_{r \in \{r_0,r_1,r_2\}} r.\pi$ after $n$ complete cycles, starting
from $P(0) = 0$.

\begin{theorem}[Linear paradox growth]
\label{thm:paradox}
$P(n) = 4n$ for all $n \geq 0$.
\end{theorem}

\begin{proof}
Count engagement events per cycle.
\texttt{ENGAGR \%r0} contributes 1 event ($r_0.\pi \mathrel{+}= 1$).
\texttt{FSPLIT \%r0 \%r1 \%r2} contributes 3 events (one per register).
\texttt{FFUSE} contributes 0 events.
The per-cycle increment is therefore $1 + 3 + 0 = 4$.  Since cycles are
identical, $P(n) = 4n$ by induction.
\end{proof}

\begin{corollary}
The three registers accumulate paradox counts in the ratio
$r_0.\pi : r_1.\pi : r_2.\pi = 2 : 1 : 1$ (both ENGAGR and FSPLIT fire on
$r_0$; only FSPLIT fires on $r_1$ and $r_2$).
\end{corollary}

\section{Empirical Verification}

The program \texttt{para\_loop.py} runs $\mathcal{K}$ on a \texttt{ParaEngine}
instance, reporting periodic snapshots.  Two checkpoints were recorded during
a single continuous session:

\begin{center}
\begin{tabular}{lrrrl}
\toprule
Checkpoint & Steps & Cycles $n$ & Paradox $P$ & $P - 4n$ \\
\midrule
$t_1$ & 625,700,000 & 208,566,666 & 834,266,668 & $+4$ \\
$t_2$ & 18,956,065,000 & 6,318,688,333 & 25,274,753,333 & $+1$ \\
\bottomrule
\end{tabular}
\end{center}

The residuals $P - 4n \in \{0, 1, 2, 3, 4\}$ arise because snapshots are
taken mid-cycle.  At $t_1$, residual $+4$ indicates capture immediately
after \texttt{FSPLIT} (all four events of cycle $n+1$ already logged); at
$t_2$, residual $+1$ indicates capture after \texttt{ENGAGR} alone.  The
residual is therefore a cycle-phase indicator, not a deviation from
\cref{thm:paradox}.

In both snapshots: active registers $= 3$, fixed registers $= 0$, Belnap
distribution $= (N{=}0,\, T{=}0,\, F{=}0,\, B{=}3)$.  The machine ran from
the first checkpoint to the second without interruption, a wall-clock
duration of approximately 13.5 hours.

\section{Remarks}

\paragraph{Non-triviality.}
The result is not the observation that ENGAGR forces $B$ and $B$ persists;
that would follow from any constant-forcing opcode.  The non-trivial point is
that FFUSE---whose semantic is the Belnap join, a \emph{non-trivial} binary
operation on $\mathbf{B}_4$---preserves $B$ specifically because $B$ is the
top element of the information order and $\vee$ is idempotent there.  A
different join (e.g.\ one that collapses $B \vee B$ to $T$) would break the
fixed point.  The Belnap lattice is the minimal structure in which paradox is
both non-explosive and self-sustaining.

\paragraph{Paraconsistency vs.\ classical contradiction.}
In a classical machine, a register holding both True and False would represent
an error state and computation would halt or diverge.  Here contradiction is
stabilized in place: the engine continues executing, paradox count grows as a
clean linear function, and no classical truth value is ever witnessed.  The
machine does not \emph{resolve} the contradiction; it \emph{runs on it}.

\section{Resolution: IFIX Cannot Collapse the Loop}

We answer the open question stated above.  There are two independent reasons
IFIX cannot drive the system to all-$T$.

\begin{theorem}[IFIX absorption by B]
\label{thm:ifix}
Let $\mathcal{K}'$ be any circular kernel that retains ENGAGR and FSPLIT
over the registers feeding into FFUSE.  Then no finite sequence of IFIX
instructions in $\mathcal{K}'$ can produce a state in which all registers
hold belief $T$ at the end of any cycle.
\end{theorem}

\begin{proof}
Two cases cover all placements of IFIX in $\mathcal{K}'$.

\textbf{Case A: IFIX fires on a register $r$ before FSPLIT touches $r$ in
the same cycle.}
FSPLIT calls \texttt{engage()} on every argument register.
\texttt{engage()} sets \texttt{flux $\leftarrow$ 11} without consulting the
\texttt{is\_fixed} property; it also updates belief to $B$ via
\texttt{\_set\_belief}.  Therefore the IFIX-$T$ state is overridden by the
subsequent FSPLIT, and $r$ returns to $B$ before FFUSE fires.

\textbf{Case B: IFIX fires on a register $r$ after FSPLIT and before
FFUSE, so $r$ enters FFUSE with belief $T$.}
At least one other FFUSE input is a register touched by FSPLIT this cycle
(by assumption that FSPLIT feeds FFUSE); that register holds $B$.
By the Belnap join: $T \vee B = B$.
Therefore FFUSE outputs $B$, not $T$.

In both cases the FFUSE output that feeds the root $r_0$ is $B$.
ENGAGR at the start of the next cycle forces $r_0$ back to $B$ regardless.
By induction over cycles, the system never reaches all-$T$.
\end{proof}

\begin{remark}[Register-level paraconsistency in K3]
A third phenomenon appears when IFIX fires on $r_0$ \emph{before} ENGAGR in
the same cycle (\textbf{K3} below).  IFIX sets \texttt{r.value = FIXED} and
belief $\leftarrow T$.  ENGAGR then sets \texttt{flux = 11} and belief
$\leftarrow B$, but does \emph{not} clear \texttt{r.value}.  The register
therefore simultaneously satisfies the linear type constraint
(\texttt{value = FIXED}, installed by IFIX) and the paradox belief state
(\texttt{flux = 11}, installed by ENGAGR).  This is paraconsistency at the
register level: the FIXED constraint and the B-state coexist without
contradiction.  The machine does not resolve the tension; it carries both
indefinitely.
\end{remark}

\subsection*{Experimental confirmation}

Three extended kernels were executed for $5 \times 10^5$ cycles each
(script \texttt{ifix\_experiment.py}):

\begin{center}
\begin{tabular}{llll}
\toprule
Kernel & IFIX placement & Belnap at cycle end & $r_0$ belief \\
\midrule
K1 & on $r_1$, after FSPLIT & $N{=}0\ T{=}1\ F{=}0\ B{=}2$ & $B$ \\
K2 & on $r_2$, after FSPLIT & $N{=}0\ T{=}1\ F{=}0\ B{=}2$ & $B$ \\
K3 & on $r_0$, before ENGAGR & $N{=}0\ T{=}0\ F{=}0\ B{=}3$ & $B$ [FIXED] \\
\bottomrule
\end{tabular}
\end{center}

In K1 and K2 the IFIXed register ends each cycle at $T$ (IFIX fires last on
it), but $r_0$ remains $B$ in all cases via Case B of \cref{thm:ifix}.
In K3 the FIXED marker persists on $r_0$ alongside flux $= 11$ and belief
$= B$, confirming the register-level paraconsistency described in the Remark
above.

\begin{corollary}
The Frobenius loop is \emph{IFIX-stable}: no circular extension of
$\mathcal{K}$ that preserves ENGAGR and FSPLIT can collapse the system to
all-$T$.  The paraconsistent fixed point is robust to the linear type
constraint.
\end{corollary}

\end{document}
